\section{Animaciones y visualizaciones}

\subsection{Animación del encuentro cercano con vectores}

Animar la trayectoria de Apophis cerca de la Tierra mostrando en cada cuadro los vectores $\vec{h}$, $\vec{e}$, $\vec{r}$ y $\vec{v}$. Observar cómo $\vec{e}$ y $\vec{h}$ cambian durante el encuentro, puesto que ya no se trata de un sistema puramente de 2 cuerpos.

\subsection{Animación del sistema solar con Apophis resaltado}

Construir una animación \textit{top-down} del sistema solar interior durante los años 2028--2030, mostrando las órbitas de Mercurio, Venus, Tierra, Marte y Apophis. Resaltar el momento del encuentro con un acercamiento automático.

\subsection{Trayectoria geocéntrica animada}

Animar la trayectoria de Apophis en un sistema de referencia centrado en la Tierra, como si se observara desde nuestro planeta. Mostrar como referencia el anillo de la órbita geoestacionaria ($42{,}164\,\mathrm{km}$) y la distancia media a la Luna ($384{,}400\,\mathrm{km}$).

\subsection{Hodógrafo animado}

Animar el hodógrafo (trayectoria del vector velocidad en el espacio de velocidades) de Apophis a lo largo de su órbita completa. Mostrar cómo se deforma la circunferencia del hodógrafo durante el encuentro con la Tierra y compararlo con el hodógrafo kepleriano puro.